El predicado a corresponde a una demostración por inducción estructural, ya que podriamos reescribirlo de la siguiente forma:

$$\paraTodo{c}{\circuito}{\alternado \composition \alternado \; c = \textito{id} \; c}$$

Para empezar, podemos utilizar las distintas definciones que el enunciado menciona para reducir esta expresión.

El lado izquierdo del predicado se reduce a:

\begin{equation*}
    \begin{aligned}
        \igualesPor{\alternado \composition \alternado \; c}{\alternado \; (\alternado \; c)}{por \formula{C}}
    \end{aligned}
\end{equation*}

El lado derecho del predicado pasa a ser simplemente $c$, pues:

\begin{equation*}
    \begin{aligned}
        \igualesPor{\textito{id} \; c}{c}{por \formula{I}}
    \end{aligned}
\end{equation*}

Entonces, lo que queremos demostrar se reduce a:

$$\paraTodo{c}{\circuito}{\alternado \; (\alternado \; c) = c}$$

Para continuar reduciendo la expresión, debemos usar todos los posibles constructores de \circuito. De esa forma, vamos a tener tres nuevos predicados, que son:

\begin{equation*}
    \paraTodo{c}{\caja}{\alternado \; (\alternado \; (\caja \; c)) = \caja \; c}
\end{equation*}

\begin{equation*}
    \paraTodo{c1, c2}{\circuito}{\alternado \; (\alternado \; (\serie{c1}{c2})) = \serie{c1}{c2}}
\end{equation*}

\begin{equation*}
    \paraTodo{ca1, ca2}{\caja}{\paraTodo{ci1, ci2}{\circuito}{\alternado \; (\alternado \; (\paralelo{ca1}{ci1}{ci2}{ca2})) = \paralelo{ca1}{ci1}{ci2}{ca2}}}
\end{equation*}

De ahora en adelante, vamos a dejar de escribir los cuantificadores por comodidad, pero de manera implicita siempre estan ahi.

Ahora veamos caso por caso.

\subsubsection*{Caso base}

Queremos ver que:

\begin{equation*}
    \paraTodo{c}{\caja}{\alternado \; (\alternado \; (\caja \; c)) = \caja \; c}
\end{equation*}

Vamos a trabajar con el lado izquierdo de la expresión.

\begin{equation*}
    \begin{aligned}
        \igualesPor{\alternado \; (\alternado \; (\caja \; c))}{\alternado \; (\caja \; (\cajaAlternada \; c))}{por \formula{AC}}
    \end{aligned}
\end{equation*}

Ahora bien, si vemos los contructores de $\caja$ tenemos que puede ser tanto $\nada$ como una $\bombilla{b}$, con algun $b$ de tipo $\bool$. Con esto en mente, volvemos a separar en casos.

Si suponemos que $c = \nada$, entonces nos queda:

\begin{equation*}
    \begin{aligned}
        \igualesPor{\alternado \; (\caja \; (\cajaAlternada \; \nada))}{\alternado \; (\caja \; \nada)}{por \formula{CAN}} \\
        \igualesPor{}{\caja \; (\cajaAlternada \; \nada)}{por \formula{AC}} \\
        \igualesPor{}{\caja \; \nada}{por \formula{CAN}} \\
    \end{aligned}
\end{equation*}

Si retomamos el lado derecho de la expresión reemplazando $c$ por $Nada$, nos queda:

$$\caja \; \nada = \caja \; \nada$$

Que es, evidentemente, verdadero.

Veamos los casos de $\bombilla{b}$.

\begin{equation*}
    \begin{aligned}
        \igualesPor{\alternado \; (\caja \; (\cajaAlternada \; \bombilla{b}))}{\alternado \; (\caja \; (\bombilla{(\notfunc \; b)}))}{por \formula{CAN}} \\
    \end{aligned}
\end{equation*}

Y al otro lado de la expresión tenemos:

$$\caja \; \bombilla{b}$$

Para resolver ese $\notfunc$ debemos volver a separar en casos (les juro que es la ultima vez), y ver que pasa con $b$. Como $b$ es de tipo $\bool$ sabemos que puede ser tanto $\true$ como $\false$. Otra vez, probemos uno por uno. Supongamos, primero, que $b = \true$.

\begin{equation*}
    \begin{aligned}
        \igualesPor{\alternado \; (\caja \; (\bombilla{\notfunc \; \true}))}{\alternado \; (\caja \; \bombilla{\false})}{por \formula{NT}} \\
        \igualesPor{}{\cajaAlternada \; (\caja \; \bombilla{\false})}{por \formula{AC}} \\
        \igualesPor{}{\caja \; (\bombilla{\notfunc \; \false})}{por \formula{CAB}} \\
        \igualesPor{}{\caja \; \bombilla{\true}}{por \formula{NF}} \\
    \end{aligned}
\end{equation*}

Esto, visiblemente, es lo mismo que teniamos del otro lado de la expresión.

El caso de $b = \false$ es analogo, como se puede ver a continuación:

\begin{equation*}
    \begin{aligned}
        \igualesPor{\alternado \; (\caja \; (\bombilla{\notfunc \; \false}))}{\alternado \; (\caja \; \bombilla{\true})}{por \formula{NF}} \\
        \igualesPor{}{\cajaAlternada \; (\caja \; \bombilla{\true})}{por \formula{AC}} \\
        \igualesPor{}{\caja \; (\bombilla{\notfunc \; \true})}{por \formula{CAB}} \\
        \igualesPor{}{\caja \; \bombilla{\false}}{por \formula{NT}} \\
    \end{aligned}
\end{equation*}

Con todo esto, el caso base de $\caja$ queda demostrado.

\subsubsection*{Paso inductivo}

